% Pass 3 cuts (level of detail against Dablander's DPhil thesis) from Chapter 2 (Background).
% Compared with his Chapter 2 (pp. 8-43, representations only), Section 3.2 (pp. 48-51, activity cliffs) and Sections 3.3.5-3.3.6 (pp. 61-64, metrics and models).
% His background covers representations in more depth than these notes plan, so nothing was cut there; his Chapter 2 ends with a table of each method's strengths and weaknesses, so one line asks for that table.
% He gives models no background at all: random forests, k-nearest neighbours and MLPs are named in one paragraph of methods, hyperparameters in tables. Metrics and splits are defined in the chapter that uses them.
% So the cuts are small: one model-section bullet that repeats the dimensionality point, software-paper detail of NoiseInject (moved to an Additional file by a note), and dataset-specific split detail (moved to Appendix "Datasets and curation").

% ===== FROM: Support vector machines notes =====
% - With the RBF kernel every feature enters every distance, so uninformative features blur it.

% ===== FROM: Evaluation methodology notes, Paragraph 2 and Paragraph 3 =====
% - Chapter 3 and Chapters 5 and 6 build scaffold groups differently (stereochemistry, acyclic molecules, group order).
% - Non-molecule datasets use their own groups: antibody sequence-identity clusters, reaction additives, chromosomes, tissue source sites.
